% GENERATED FILE. Edit canonical lessons, then run npm run generate:books.
% canonical-source-hash: fnv1a64:252efb2d7cb7c8c3
% canonical-lessons: ML-C387-salkkaram, ML-C387-virunnukaran, ML-C387-catannu, ML-C387-aniyuka, ML-C387-orukkuka

\chapter{The Feast and the Ceremony}
\label{ch:ml-the-feast-and-the-ceremony}
\hlchaptermodality{387}

\begin{chapteropening}
\textbf{By the end of this chapter:} \emph{I can say a banquet, a house guest, a ceremony, to put on finery and to make ready.}

Complete the last lesson of chapter 387: I can say a banquet, a house guest, a ceremony, to put on finery and to make ready.
\end{chapteropening}

\section[\texorpdfstring{\ml{സൽക്കാരം} (salkkāraṁ)}{സൽക്കാരം (salkkāraṁ)}]{\ml{സൽക്കാരം} (salkkāraṁ) — a banquet}

\label{lesson:ML-C387-salkkaram}



\begin{quote}
Before the new one: say the Malayalam for to invite, then the Malayalam for to receive.
\end{quote}

\subsection*{You'll want to know: \ml{സൽക്കാരം}}
\textbf{\ml{സൽക്കാരം}} — \emph{salkkāraṁ} — \textquotedblleft{}a banquet\textquotedblright{}.

\begin{quote}
\ml{മാനേജർ} \ml{ജീവനക്കാർക്ക്} \ml{ഒരു} \ml{സൽക്കാരം} \ml{നടത്തി}.
\end{quote}

\begin{quote}
\emph{mānējar jīvanakkārkkŭ oru salkkāraṁ naṭatti.}
\end{quote}

\begin{quote}
The manager held a banquet for the staff.
\end{quote}

A banquet is held: \ml{സൽക്കാരം} \ml{നടത്തുക}.

\begin{cousinweb}
\textbf{\ml{സൽക്കാരം}} (\emph{salkkāraṁ}) is borrowed from Sanskrit \emph{satkāra}, \textquotedblleft{}kind treatment, hospitality\textquotedblright{} (CDIAL 13108), from \emph{sat}, \textquotedblleft{}good\textquotedblright{}, and \emph{kāra}, \textquotedblleft{}a doing\textquotedblright{}.
\end{cousinweb}

\begin{grammarlens}[title={What you've built}]
One more word for telling what happened.
\end{grammarlens}

\subsection*{Your turn}
\begin{itemize}
\raggedright
  \item \emph{Say it:} \emph{salkkāraṁ}
  \item \emph{Say it:} \emph{salkkāraṁ}, once more
  \item \emph{Say it:} \emph{mānējar jīvanakkārkkŭ oru salkkāraṁ naṭatti.} — The manager held a banquet for the staff.
  \item \emph{Recall:} say the Malayalam for to invite, then the Malayalam for to receive, then say \emph{salkkāraṁ} again
\end{itemize}

\begin{tcolorbox}[breakable,colback=teal!4,colframe=teal!35!black,title={Before you move on}]
What does \ml{സൽക്കാരം} mean? (\textquotedblleft{}A banquet\textquotedblright{}.) Say it once more.
\end{tcolorbox}
\section[\texorpdfstring{\ml{വിരുന്നുകാരൻ} (virunnukāran)}{വിരുന്നുകാരൻ (virunnukāran)}]{\ml{വിരുന്നുകാരൻ} (virunnukāran) — a house guest; a visitor}

\label{lesson:ML-C387-virunnukaran}



\begin{quote}
Before the new one: say the Malayalam for to receive, then the Malayalam for a banquet.
\end{quote}

\subsection*{You'll want to know: \ml{വിരുന്നുകാരൻ}}
\textbf{\ml{വിരുന്നുകാരൻ}} — \emph{virunnukāran} — \textquotedblleft{}a house guest; a visitor\textquotedblright{}.

\begin{quote}
\ml{അവധിക്ക്} \ml{വീട്ടിൽ} \ml{ഒരുപാട്} \ml{വിരുന്നുകാർ} \ml{വന്നു}.
\end{quote}

\begin{quote}
\emph{avadhikkŭ vīṭṭil orupāṭŭ virunnukār vannu.}
\end{quote}

\begin{quote}
During the holidays many visitors came to the house.
\end{quote}

Plural: \ml{വിരുന്നുകാർ}, \emph{virunnukār}.

From \ml{വിരുന്ന്}, a feast: the guest who comes to stay and eat.

\begin{cousinweb}
\textbf{\ml{വിരുന്നുകാരൻ}} (\emph{virunnukāran}) is native Dravidian: \ml{വിരുന്ന്}, \textquotedblleft{}a feast\textquotedblright{}, which DEDR 5415 sets beside Tamil \emph{viruntu}, \textquotedblleft{}a feast, a guest\textquotedblright{}, with the ending -\ml{കാരൻ}.
\end{cousinweb}

\begin{grammarlens}[title={What you've built}]
One more word for telling what happened.
\end{grammarlens}

\subsection*{Your turn}
\begin{itemize}
\raggedright
  \item \emph{Say it:} \emph{virunnukāran}
  \item \emph{Say it:} \emph{virunnukāran}, once more
  \item \emph{Say it:} \emph{avadhikkŭ vīṭṭil orupāṭŭ virunnukār vannu.} — During the holidays many visitors came to the house.
  \item \emph{Recall:} say the Malayalam for to receive, then the Malayalam for a banquet, then say \emph{virunnukāran} again
\end{itemize}

\begin{tcolorbox}[breakable,colback=teal!4,colframe=teal!35!black,title={Before you move on}]
What does \ml{വിരുന്നുകാരൻ} mean? (\textquotedblleft{}A house guest; a visitor\textquotedblright{}.) Say it once more.
\end{tcolorbox}
\section[\texorpdfstring{\ml{ചടങ്ങ്} (caṭaṅṅŭ)}{ചടങ്ങ് (caṭaṅṅŭ)}]{\ml{ചടങ്ങ്} (caṭaṅṅŭ) — a ceremony; a function}

\label{lesson:ML-C387-catannu}



\begin{quote}
Before the new one: say the Malayalam for a banquet, then the Malayalam for a house guest.
\end{quote}

\subsection*{You'll want to know: \ml{ചടങ്ങ്}}
\textbf{\ml{ചടങ്ങ്}} — \emph{caṭaṅṅŭ} — \textquotedblleft{}a ceremony; a function\textquotedblright{}.

\begin{quote}
\ml{വിവാഹത്തിന്റെ} \ml{ചടങ്ങ്} \ml{രാവിലെ} \ml{പത്ത്} \ml{മണിക്ക്} \ml{തുടങ്ങി}.
\end{quote}

\begin{quote}
\emph{vivāhattinṟe caṭaṅṅŭ rāvile pattŭ maṇikkŭ tuṭaṅṅi.}
\end{quote}

\begin{quote}
The wedding ceremony began at ten in the morning.
\end{quote}

Also a mere formality: \ml{അത്} \ml{ഒരു} \ml{ചടങ്ങ്} \ml{മാത്രം}, it is only a formality.

\begin{cousinweb}
\textbf{\ml{ചടങ്ങ്}} (\emph{caṭaṅṅŭ}) is a native Malayalam word whose origin is not settled. It names any ceremony, from a wedding to the opening of a shop.
\end{cousinweb}

\begin{grammarlens}[title={What you've built}]
One more word for telling what happened.
\end{grammarlens}

\subsection*{Your turn}
\begin{itemize}
\raggedright
  \item \emph{Say it:} \emph{caṭaṅṅŭ}
  \item \emph{Say it:} \emph{caṭaṅṅŭ}, once more
  \item \emph{Say it:} \emph{vivāhattinṟe caṭaṅṅŭ rāvile pattŭ maṇikkŭ tuṭaṅṅi.} — The wedding ceremony began at ten in the morning.
  \item \emph{Recall:} say the Malayalam for a banquet, then the Malayalam for a house guest, then say \emph{caṭaṅṅŭ} again
\end{itemize}

\begin{tcolorbox}[breakable,colback=teal!4,colframe=teal!35!black,title={Before you move on}]
What does \ml{ചടങ്ങ്} mean? (\textquotedblleft{}A ceremony; a function\textquotedblright{}.) Say it once more.
\end{tcolorbox}
\section[\texorpdfstring{\ml{അണിയുക} (aṇiyuka)}{അണിയുക (aṇiyuka)}]{\ml{അണിയുക} (aṇiyuka) — to put on finery; to adorn oneself}

\label{lesson:ML-C387-aniyuka}



\begin{quote}
Before the new one: say the Malayalam for a house guest, then the Malayalam for a ceremony.
\end{quote}

\subsection*{You'll want to know: \ml{അണിയുക}}
\textbf{\ml{അണിയുക}} — \emph{aṇiyuka} — \textquotedblleft{}to put on finery; to adorn oneself\textquotedblright{}.

\begin{quote}
\ml{വധു} \ml{ചുവന്ന} \ml{സാരി} \ml{അണിഞ്ഞു}.
\end{quote}

\begin{quote}
\emph{vadhu cuvanna sāri aṇiññu.}
\end{quote}

\begin{quote}
The bride put on a red sari.
\end{quote}

Past: \ml{അണിഞ്ഞു}, \emph{aṇiññu}.

Of jewels and fine clothes; \ml{ധരിക്കുക} is the everyday word.

\begin{cousinweb}
\textbf{\ml{അണിയുക}} (\emph{aṇiyuka}) is native Dravidian: DEDR 116 lists Malayalam \emph{aṇiyuka}, \textquotedblleft{}to wear (jewels, arms), to decorate oneself\textquotedblright{}, beside Tamil \emph{aṇi}, \textquotedblleft{}an ornament\textquotedblright{}.
\end{cousinweb}

\begin{grammarlens}[title={What you've built}]
One more word for telling what happened.
\end{grammarlens}

\subsection*{Your turn}
\begin{itemize}
\raggedright
  \item \emph{Say it:} \emph{aṇiyuka}
  \item \emph{Say it:} \emph{aṇiyuka}, once more
  \item \emph{Say it:} \emph{vadhu cuvanna sāri aṇiññu.} — The bride put on a red sari.
  \item \emph{Recall:} say the Malayalam for a house guest, then the Malayalam for a ceremony, then say \emph{aṇiyuka} again
\end{itemize}

\begin{tcolorbox}[breakable,colback=teal!4,colframe=teal!35!black,title={Before you move on}]
What does \ml{അണിയുക} mean? (\textquotedblleft{}To put on finery; to adorn oneself\textquotedblright{}.) Say it once more.
\end{tcolorbox}
\section[\texorpdfstring{\ml{ഒരുക്കുക} (orukkuka)}{ഒരുക്കുക (orukkuka)}]{\ml{ഒരുക്കുക} (orukkuka) — to make ready; to arrange}

\label{lesson:ML-C387-orukkuka}



\begin{quote}
Before the new one: say the Malayalam for a ceremony, then the Malayalam for to put on finery.
\end{quote}

\subsection*{You'll want to know: \ml{ഒരുക്കുക}}
\textbf{\ml{ഒരുക്കുക}} — \emph{orukkuka} — \textquotedblleft{}to make ready; to arrange\textquotedblright{}.

\begin{quote}
\ml{അവർ} \ml{വലിയ} \ml{സദ്യ} \ml{ഒരുക്കി}.
\end{quote}

\begin{quote}
\emph{avar valiya sadya orukki.}
\end{quote}

\begin{quote}
They got a big feast ready.
\end{quote}

Past: \ml{ഒരുക്കി}, \emph{orukki}.

The partner of \ml{ഒരുങ്ങുക}, to get ready: you get yourself ready, you get a thing ready.

\begin{cousinweb}
\textbf{\ml{ഒരുക്കുക}} (\emph{orukkuka}) is native Dravidian: DEDR 992 lists Malayalam \emph{orukkuka}, \textquotedblleft{}to prepare, to get ready\textquotedblright{}, beside \ml{ഒരുങ്ങുക}, \textquotedblleft{}to be ready\textquotedblright{}.
\end{cousinweb}

\begin{grammarlens}[title={What you've built}]
One more word for telling what happened.
\end{grammarlens}

\subsection*{Your turn}
\begin{itemize}
\raggedright
  \item \emph{Say it:} \emph{orukkuka}
  \item \emph{Say it:} \emph{orukkuka}, once more
  \item \emph{Say it:} \emph{avar valiya sadya orukki.} — They got a big feast ready.
  \item \emph{Recall:} say the Malayalam for a ceremony, then the Malayalam for to put on finery, then say \emph{orukkuka} again
\end{itemize}

\begin{tcolorbox}[breakable,colback=teal!4,colframe=teal!35!black,title={Before you move on}]
What does \ml{ഒരുക്കുക} mean? (\textquotedblleft{}To make ready; to arrange\textquotedblright{}.) Say it once more.
\end{tcolorbox}
